\documentclass[11pt]{article}
\usepackage[final]{acl}
\usepackage{times}
\usepackage{latexsym}
\usepackage[T1]{fontenc}
\usepackage[utf8]{inputenc}
\usepackage{microtype}
\usepackage{graphicx}
\usepackage{booktabs}
\usepackage{amssymb}
\usepackage{array}
\usepackage{calc}   % pandoc's p{} column specs use calc's \real for proportional widths
\providecommand{\tightlist}{\setlength{\itemsep}{0pt}\setlength{\parskip}{0pt}}
\DeclareUnicodeCharacter{00A7}{\S}
\DeclareUnicodeCharacter{00B1}{$\pm$}
\DeclareUnicodeCharacter{00D7}{$\times$}
\DeclareUnicodeCharacter{00F7}{$\div$}
\DeclareUnicodeCharacter{0394}{$\Delta$}
\DeclareUnicodeCharacter{03B2}{$\beta$}
\DeclareUnicodeCharacter{03B8}{$\theta$}
\DeclareUnicodeCharacter{03BB}{$\lambda$}
\DeclareUnicodeCharacter{03C1}{$\rho$}
\DeclareUnicodeCharacter{2013}{--}
\DeclareUnicodeCharacter{2014}{---}
\DeclareUnicodeCharacter{2016}{$\|$}
\DeclareUnicodeCharacter{2026}{\ldots}
\DeclareUnicodeCharacter{2032}{$'$}
\DeclareUnicodeCharacter{2081}{$_1$}
\DeclareUnicodeCharacter{2192}{$\rightarrow$}
\DeclareUnicodeCharacter{2212}{$-$}
\DeclareUnicodeCharacter{2248}{$\approx$}
\DeclareUnicodeCharacter{2264}{$\leq$}
\DeclareUnicodeCharacter{2265}{$\geq$}
% Report the geometry the template actually produces, rather than the geometry its
% source appears to declare. acl.sty's \def\@normalsize is LaTeX 2.09 legacy and is
% never invoked by LaTeX2e, so the body font comes from \documentclass[11pt] instead.
\makeatletter
\AtBeginDocument{\typeout{ACLDIAG fsize=\f@size\space
  baselineskip=\the\baselineskip\space textheight=\the\textheight\space
  columnwidth=\the\columnwidth\space textwidth=\the\textwidth}}
\makeatother
\title{Can Critical Periods in Language Models Be Reopened?}
\author{}
\begin{document}
\maketitle
\begin{abstract}
Critical periods --- windows after which a skill becomes hard to acquire regardless of
further data --- are usually studied in models whose pretraining history cannot be
controlled. We train transformers of 11M--60M parameters from random initialization on a
\textasciitilde10M-word corpus, injecting synthetic arithmetic and long-distance \emph{wh}-licensing on
delay schedules that hold total skill exposure fixed and vary only \emph{when} that exposure
begins.

Delaying arithmetic exposure to the midpoint of training produces a persistent
endpoint deficit against seed-matched early-exposure anchors --- 0.214 ± 0.055 at 23M and
0.229 ± 0.037 at 60M, in the same direction at every seed. Long-distance \emph{wh}-licensing
shows no detectable deficit under five distinct probes, inverting our pre-registered
prediction that syntax would show the \emph{stronger} critical period. We claim this
dissociation in absolute terms and not as differential susceptibility, because the two
skills are not equally observable: the syntax anchor sits 0.028 above its metric's
chance floor against arithmetic's 0.791, so a proportionally matched deficit would have
fallen below the design's resolution (MDE 0.048). Four of the five probes ask whether a
critical period exists \emph{given the acquisition that occurred}; only a representational
probe --- which carries independent dynamic range (selectivity 0.60, p \textless{} 1e-8) and is flat
across timing (p = 0.48) --- asks whether enough was acquired for one to have been
detectable at all. That probe, not the count of nulls, is what the dissociation rests on.

Where the deficit exists, it reopens. Shifting the learning-rate schedule at the delay
onset restores endpoint accuracy to 0.767 ± 0.014 at 23M and 0.796 ± 0.022 at 60M,
against untreated nulls of 0.577 ± 0.051 and 0.600 ± 0.040 --- at 23M with no measurable
cost to held-out perplexity, at 60M with a cost bounded above by 1.67 ppl that is
plausibly an annealing-horizon artifact. Both recovered endpoints fall inside a
pre-registered band, so recovery converges on a scale-invariant level rather than
tracking capacity, and it collapses the cross-seed variance the untreated deficit shows.
Concentrated re-exposure to the missed skill does not reproduce this: a data-only cram
phase stays near the null, while cram combined with the schedule shift is
indistinguishable from the shift alone. Closure at this onset is therefore largely a
property of the optimizer's decayed learning-rate schedule rather than of entrenched
representations.

The two results are timing effects at different levels: \emph{which} skills a late start
damages is skill-specific, while the damage itself is carried by the optimizer's
schedule rather than by representations that have hardened against revision.

\end{abstract}
\section{§1 Introduction}\label{introduction}

A critical period is a window during which experience shapes a system, after which the same
experience has much less effect. Whether anything like it exists in language models is
usually asked of models whose training history cannot be inspected: a released checkpoint
has already seen whatever the web contained, in whatever order, so the timing of any
particular skill's arrival is unknown and unmanipulable. Fine-tuning such a model tells us
what it can still learn, not what it could have learned had the material arrived earlier.

We take the other route. Every model here is trained from random initialization on a fixed
\textasciitilde10M-word corpus, with a target skill injected into the training stream on a schedule we
control. Conditions differ in \emph{when} the skill arrives and in nothing else --- total skill
tokens are held constant --- so a difference in final performance is a timing effect, not a
dose effect. This costs us scale, the largest model being 60.3M parameters, and buys the one
thing a pretrained checkpoint cannot provide: knowing when the model first saw arithmetic.

Within that design we find, first, that \textbf{delaying arithmetic exposure to the midpoint of
training leaves a deficit that the remaining half of training does not close}, at both
confirmatory sizes and at every seed. Second, that \textbf{long-distance \emph{wh}-licensing shows no
such deficit} under five distinct probes --- inverting our pre-registered prediction that
syntax, the more structurally scaffolded skill, would be the more fragile one. We state that
dissociation in absolute terms rather than as differential susceptibility; §3.3 gives the
measurement reason. Third, that \textbf{the deficit is reopenable}: restoring the learning-rate
schedule at the point of delayed exposure recovers the endpoint to within 90 \% of an
early-exposure anchor at both sizes, and to nearly the same absolute level at each. Fourth,
that \textbf{concentrated re-exposure does not reproduce this} --- extra skill density on top of a
restored schedule adds nothing, while a restored schedule without it recovers the skill.
What closes the window at this onset is therefore largely the optimizer's decayed
learning-rate schedule rather than representations hardened against revision.

Two things follow for how we report. A delayed model's endpoint is strongly seed-dependent
while its early-exposure anchor is not, so we report absolute endpoints beside their nulls
rather than recovery as a percentage of a gap whose denominator moves. And because several
results here are nulls, we read non-significance off the minimum detectable effect rather
than the \emph{p}-value --- two structurally identical size comparisons in this paper carry
opposite evidential weight for that reason (§3.1, §5.3).

We registered three predictions before running these experiments: that syntax would show the
stronger critical period, that embedding reset would be the most effective intervention, and
that a reset-driven recovery would build new circuitry rather than repair the original. §7.2
states what happened to each. A reader holding the pre-registration will find all three
answered there, including the one that could not be tested.

\begin{center}\rule{0.5\linewidth}{0.5pt}\end{center}

\section{§2 Methods}\label{methods}

\subsection{2.1 Models and corpus}\label{models-and-corpus}

All models are decoder-only transformers trained from random initialization; \textbf{no pretrained
checkpoint is used anywhere in this work, including for baselines.} This is the design
constraint that makes the rest of the paper meaningful --- a model that has already passed
through an uncontrolled pretraining history cannot be asked when it first met a skill.

The base corpus is the BabyLM ``strict-small'' collection (\textasciitilde10M words), tokenized with a
byte-pair encoder trained on that corpus, vocabulary 16,384. Three configurations are used:

\begin{table}[t]
\centering\small
\begin{tabular}{@{}lrrrrr@{}}
\toprule\noalign{}
tag & layers & heads & width & total params & non-embedding params \\
\midrule\noalign{}

11M & 4 & 8 & 256 & 11.54M & 3.15M \\
23M & 6 & 8 & 384 & 23.20M & 10.62M \\
60M & 8 & 10 & 640 & 60.30M & 39.33M \\
\bottomrule
\end{tabular}
\end{table}


Both columns are given because untied embeddings make them tell different stories. At
vocabulary 16,384 the input embedding and output projection are 4.19M parameters each at the
smallest size --- \textbf{72.7 \% of that model, against 34.8 \% at the largest} --- so the three
configurations span 5.2× in total parameters but 12.5× in the parameters that do computation.
Size comparisons in this paper are stated in total parameters, which is the axis our
pre-registered analyses used; readers translating our scale results into non-embedding terms
should note that the effective range is wider than the labels suggest, and that the smallest
model is unusually embedding-dominated (§2.5).

Context length is 512 tokens. Positional information is rotary (θ = 10,000); normalization is
pre-norm; activations are GELU. \textbf{Input embeddings and the output projection are untied} ---
a deliberate choice, since one of our interventions resets the embedding layer and tied
weights would confound an input-representation reset with a change to the prediction head.
Attention is \texttt{F.scaled\_dot\_product\_attention}, which dispatches to FlashAttention-2 kernels
on our GPU. Training uses bf16 autocast, micro-batches of 8 × 512 tokens with 4 gradient
accumulation steps, AdamW (β = 0.9/0.95, weight decay 0.1, gradient clipping 1.0), a peak
learning rate of 6e-4 with 300 warmup steps and cosine decay to 0.1 × peak. \texttt{torch.compile}
is available in the trainer but was \textbf{disabled for the runs reported here}. Every run
reported here is 5,000 optimizer steps over 81.92M tokens. Checkpoint spacing varies by run
family (100 steps for the 23M arithmetic runs, 1,250 at 11M and 60M, 2,500 for the size
sweep); it sets how finely a trajectory can be re-read afterwards and nothing else, since
delay and intervention triggers fire on exact step count.

\subsection{2.2 Delay schedule and exposure matching}\label{delay-schedule-and-exposure-matching}

A run's data stream is a mixture of the base corpus and a skill stream. The skill stream has
an \textbf{onset step} and a \textbf{density}, both fixed before training starts, and set so that total
skill tokens are constant across delay conditions: a delay-0 run spreads its skill exposure
across all 5,000 steps, while a delay-50 run delivers the same token count at higher density
from step 2,500 onward. Onsets are specified as fractions of total steps and converted to
exact step numbers before the run begins, so triggers fire on step count and never on
wall-clock time.

Each run is identified by a \texttt{run\_id} encoding skill, delay, exposure fraction, size and seed;
checkpoints store optimizer state and a metadata dictionary alongside weights. \textbf{A \texttt{run\_id}
is the identity of a trained model, not a label for a configuration} --- where two runs share
a configuration but were trained separately, we treat them as different models and never
average or interchange their numbers.

\subsection{2.3 Measures}\label{measures}

Three measurements recur, and each is fixed once here:

\begin{itemize}
\tightlist
\item
  \textbf{Arithmetic}: teacher-forced exact match at step 5,000 on 500 held-out problems.
\item
  \textbf{Syntax}: a continuation-scored licensing interaction on \emph{wh}-extraction minimal pairs.
\item
  \textbf{General capability}: held-out validation perplexity, always reported as a three-number
  decomposition (early anchor, delayed null, treated arm) rather than as a single delta.
\end{itemize}

The syntax metric requires a word of justification, because the obvious one fails. Raw
BLiMP-style \emph{wh}-vs-\emph{that} pairs differ by a single word, so their score is dominated by the
unigram frequency prior for ``that'': an untrained-for-syntax model scores 0.000 on one
paradigm and 1.000 on its complement while knowing nothing about gap licensing. Taking the
interaction of the with-gap and no-gap conditions cancels that main effect --- a purely
frequency-driven model lands at exactly 0.500 --- so only genuine gap-licensing knowledge moves
the score. All syntax numbers in this paper are that interaction, and 0.500 is its floor.

\subsection{2.4 Eval-set hygiene}\label{eval-set-hygiene}

Exactly one arithmetic problem set is held out of training. An earlier version of our circuit
analysis generated its evaluation items per-run, and those draws overlapped the training
problems by 17--19 \% in four of six cases; all affected results were re-scored on the held-out
set, and the correction changed a conclusion (§5.2). On identical items our two evaluation
harnesses agree to \textasciitilde0.005, so accuracies from different harnesses are comparable to about
0.01 and no finer. We report this rather than only the corrected numbers, because the
released logs contain both generations.

\subsection{2.5 The smallest size is excluded by a pre-registered floor}\label{the-smallest-size-is-excluded-by-a-pre-registered-floor}

Our pre-registration required an early-exposure anchor to reach 0.60 before a delay
comparison at that size could be interpreted. The 11M anchor reaches 0.538 and has not
plateaued, so its delay comparison is between two under-trained models and is \textbf{excluded by
the pre-registered floor as a capacity/budget boundary}. Its numbers are reported in §3.1
because they are nominally significant and a reader recomputing them will find them; they do
not support a claim in either direction about a critical period at that scale.

\subsection{2.6 Analysis designs that were planned and not run}\label{analysis-designs-that-were-planned-and-not-run}

We pre-registered three analyses: a mixed-effects model of final accuracy on delay × size, a
factorial model of recovery percentage on intervention × skill × delay, and a model of
circuit overlap on intervention × skill. None is estimable on the design that survived. After
the floor exclusion the first has twelve observations with one per seed × cell and no
within-cell replication for a random effect. The second has a banned denominator --- recovery
percentage divides by the seed-variable null --- and two of its three factors are constants.
The third has both factors constant. We replaced them with per-seed paired tests, an
absolute-endpoint gate, and a rank correlation with exact permutation \emph{p}, respectively, and
say so here rather than presenting the replacements as the original plan.

\subsection{2.7 Multiplicity}\label{multiplicity}

Confirmatory tests are Holm-corrected within a declared family and reported both uncorrected
and adjusted: the Detect family (paired deficit at two sizes), the Reopen family (lift over
the matched null at two sizes), and the agreement-slope family. Gates whose thresholds were
fixed in advance are decision rules, not tests, and carry no \emph{p}. Exploratory analyses are
flagged in-text and never carry a headline.

\begin{center}\rule{0.5\linewidth}{0.5pt}\end{center}

\section{§3 Detect}\label{detect}

\begin{figure}[t]
\centering
\rule{\linewidth}{36.00\baselineskip}
\caption{Placeholder for the §3 Detect figure.}
\end{figure}

\subsection{3.1 Delayed arithmetic exposure produces a persistent endpoint deficit}\label{delayed-arithmetic-exposure-produces-a-persistent-endpoint-deficit}

Every model in this section sees the same number of arithmetic tokens; the conditions
differ only in \emph{when} those tokens enter the stream. The delayed condition (d50) begins
arithmetic exposure at the midpoint of training, the early condition (d0) at step 0, and
both are read at step 5000 on the same 500 held-out problems. Because seeds are crossed ---
the same seed gives the same initialization in both conditions --- the design-appropriate
statistic is the per-seed paired deficit, tested separately at each size and Holm-corrected
across the two confirmatory sizes.

\textbf{Table 1.} Arithmetic exact-match accuracy at step 5000, per seed and pooled.

\begin{table}[t]
\centering\small
\begin{tabular}{@{}
  >{\raggedright\arraybackslash}p{(\linewidth - 8\tabcolsep) * \real{0.2000}}
  >{\raggedleft\arraybackslash}p{(\linewidth - 8\tabcolsep) * \real{0.2000}}
  >{\raggedleft\arraybackslash}p{(\linewidth - 8\tabcolsep) * \real{0.2000}}
  >{\raggedleft\arraybackslash}p{(\linewidth - 8\tabcolsep) * \real{0.2000}}
  >{\raggedleft\arraybackslash}p{(\linewidth - 8\tabcolsep) * \real{0.2000}}@{}}
\toprule\noalign{}
\begin{minipage}[b]{\linewidth}\raggedright
size
\end{minipage} & \begin{minipage}[b]{\linewidth}\raggedleft
seed
\end{minipage} & \begin{minipage}[b]{\linewidth}\raggedleft
early (d0)
\end{minipage} & \begin{minipage}[b]{\linewidth}\raggedleft
delayed (d50)
\end{minipage} & \begin{minipage}[b]{\linewidth}\raggedleft
deficit
\end{minipage} \\
\midrule\noalign{}

23M & 1 & 0.798 & 0.524 & 0.274 \\
& 2 & 0.792 & 0.626 & 0.166 \\
& 3 & 0.784 & 0.582 & 0.202 \\
& \textbf{mean ± SD} & \textbf{0.791 ± 0.007} & \textbf{0.577 ± 0.051} & \textbf{0.214 ± 0.055} \\
60M & 1 & 0.798 & 0.574 & 0.224 \\
& 2 & 0.840 & 0.646 & 0.194 \\
& 3 & 0.848 & 0.580 & 0.268 \\
& \textbf{mean ± SD} & \textbf{0.829 ± 0.027} & \textbf{0.600 ± 0.040} & \textbf{0.229 ± 0.037} \\
\bottomrule
\end{tabular}
\end{table}


The deficit is present at every seed at both sizes: paired \emph{t}(2) = 6.74, \emph{p} = 0.0107 at
23M and \emph{t}(2) = 10.64, \emph{p} = 0.0044 at 60M, both surviving Holm correction within the
Detect family (adjusted 0.0107 and 0.0087; Hedges \emph{g} = 2.22 and 3.51). The pre-registered
Welch gate, which required a gap of at least 0.15, is cleared at both sizes (60M one-sided
\emph{p} = 0.0010). A model that receives its arithmetic exposure late does not catch up by the
end of training, despite having received all of it.

The deficit is of \textbf{similar magnitude at both sizes}: the 60M--23M difference is 0.015,
Welch \emph{p} = 0.72, with a minimum detectable effect at 80 \% power of \textbf{0.144}. That null
constrains --- it excludes any size effect larger than about two-thirds of the deficit itself
--- but it is not a claim that size is irrelevant, and it should not be read as one. We
return to this in §5, where a structurally identical size comparison carries the opposite
evidential weight for one reason only: its MDE exceeds the entire range its outcome
occupies. The two nulls are best read as a pair.

\textbf{The severity of the deficit is seed-variable, and the variability lives on the delayed
side.} The early anchor is tight across seeds (SD 0.007 at 23M) while the delayed
endpoint is not (SD 0.051), so per-seed deficits range from 0.166 to 0.274 at 23M against a
baseline that barely moves. Seed 1 is the deepest-hit on both readings --- lowest delayed
endpoint and largest deficit --- and reproduces as such when retrained from scratch. This is
why every recovery result in §4 is reported as an absolute endpoint beside its null rather
than as a percentage of the gap: the gap is a seed-variable denominator, so a
recovery-percentage is unstable by construction even when the underlying effect is not.

\textbf{The 11M models are excluded by a pre-registered floor.} Their paired deficit is
nominally significant (0.181 ± 0.087, \emph{p} = 0.0345, \emph{g} = 1.19), and we report it because a
reader recomputing our numbers will find it. It is excluded because the 11M early anchor
(0.538 ± 0.018) never plateaued and sits below the pre-registered 0.60 floor: the
comparison is between two under-trained models, so the difference cannot be attributed to
timing rather than to capacity or budget. This is a capacity boundary, not evidence that no
critical period exists at 11M.

Intermediate onsets (d25, d55, d75) were run at 23M seed 1 only. They are plotted for
completeness without a fitted trend; at \emph{n} = 1 per onset they describe this run and support
no claim about the shape of the timing curve.

\subsection{\texorpdfstring{3.2 Long-distance \emph{wh}-licensing shows no detectable critical period}{3.2 Long-distance wh-licensing shows no detectable critical period}}\label{long-distance-wh-licensing-shows-no-detectable-critical-period}

The same design applied to long-distance \emph{wh}-licensing yields nothing. The five probes
below are \textbf{five different hypotheses about where an effect could hide, not five tests of
one hypothesis}; each carries its own positive control or power statement, and the argument
is that pairing rather than the count.

\begin{enumerate}
\def\labelenumi{\arabic{enumi}.}
\tightlist
\item
  \textbf{Endpoint by delay} (\emph{n} = 3): no gap, against a delay-0 baseline that itself ranges
  0.520--0.586 over eight seeds --- a spread exceeding any candidate effect --- on the same
  harness that resolves arithmetic's 0.214.
\item
  \textbf{Acquisition dynamics}: Welch \emph{p} = 0.46, underpowered alone and carried by probe 3.
\item
  \textbf{Rate / bimodality}: a pre-registered compound gate over \textbf{72 runs} (24 seeds × 3
  sizes) returns DISCONFIRM --- the mean rises monotonically with size (0.528 → 0.538 →
  0.552) but fails its floor at threshold 0.56 (top size 0.208 \textless{} 0.40) and its trend at
  0.57 (\emph{p} = 0.074), passing only at the permissive 0.55 --- and this powered test, sized in
  advance by simulation, is why long-distance syntax left the main grid.
\item
  \textbf{Representational probe}: delay-0 versus delay-50 \emph{p} = 0.48 \emph{while the same probe
  detects the structure at selectivity 0.60, p \textless{} 1e-8} --- the positive control passes.
\item
  \textbf{Learning-rate budget} (hilr50, \emph{n} = 3): ΔLD = +0.011 ± 0.005, with Δlocal exceeding
  ΔLD at every seed and validation perplexity improving --- a non-specific whole-model lift,
  not a rescue of the delayed skill.
\end{enumerate}

Probe 4 is the load-bearing one, for a reason that the other four cannot supply. Probes 1,
2, 3 and 5 all ask whether a critical period is visible \emph{in the behaviour that was
acquired}. Probe 4 asks whether the underlying structure is encoded at all, and finds that
it is --- and that its encoding does not depend on when exposure began. \textbf{The structure is
learned; what is timing-invariant is the readout.}

This null is scoped to ≤ 60.3M parameters, the \textasciitilde10M-word BabyLM budget, and the
continuation-scored licensing-interaction metric. Outside that scope it says nothing.

\subsection{3.3 Skill dissociation}\label{skill-dissociation}

Under one design, one pipeline and matched total exposure, delayed exposure costs
arithmetic 0.214 ± 0.055 and costs long-distance \emph{wh}-licensing nothing measurable across
five probes. That contrast is the paper's central detection-side result, and it inverts our
pre-registered prediction that syntax would show the \emph{stronger} critical period --- a
prediction we restate with its outcome as a standalone item in §7.2, since a reader holding
our pre-registration will look for it there rather than here.

\textbf{It must be stated in absolute terms, and not as differential susceptibility, because the
two skills are not equally observable.} The asymmetry is large enough to state as a table
rather than to note as a caveat:

\textbf{Table 2.} Why the dissociation is scoped.

\begin{table}[t]
\centering\small
\begin{tabular}{@{}
  >{\raggedright\arraybackslash}p{(\linewidth - 4\tabcolsep) * \real{0.3333}}
  >{\raggedleft\arraybackslash}p{(\linewidth - 4\tabcolsep) * \real{0.3333}}
  >{\raggedleft\arraybackslash}p{(\linewidth - 4\tabcolsep) * \real{0.3333}}@{}}
\toprule\noalign{}
\begin{minipage}[b]{\linewidth}\raggedright
\end{minipage} & \begin{minipage}[b]{\linewidth}\raggedleft
arithmetic
\end{minipage} & \begin{minipage}[b]{\linewidth}\raggedleft
syntax LD
\end{minipage} \\
\midrule\noalign{}

early-exposure anchor & 0.791 ± 0.007 & 0.528 (24 seeds) \\
chance floor & \textasciitilde0 (exact match) & 0.500 (by metric construction) \\
\textbf{headroom above floor} & \textbf{0.791} & \textbf{0.0283} \\
observed deficit & 0.214 (27.1 \% of headroom) & none detected \\
a proportionally matched deficit would be & --- & 0.0077 (0.49 seed-SD) \\
MDE at 80 \% power, \emph{n} = 3/group & --- & \textbf{0.0482} \\
\bottomrule
\end{tabular}
\end{table}


The long-distance design's MDE \textbf{exceeds its entire headroom by 1.7×}. A total collapse of
the skill to chance would be a change of 0.0283 --- smaller than the smallest effect the
design could resolve --- and a susceptibility proportionally identical to arithmetic's would
require \emph{n} ≈ 67 per group. The behavioural null therefore cannot, on its own, separate
\emph{syntax is less susceptible} from \emph{syntax had no room to show susceptibility}.

Two things follow, and only one of them is a rescue. The learning-rate-budget null (probe
5) does \textbf{not} relieve the problem: it establishes that the ceiling is not an
optimization-budget artifact, which makes the ceiling real and immovable by that lever, and
so \emph{preserves} the dynamic-range limit rather than lifting it. The representational probe
(probe 4) does: it is the one readout not bounded by the 0.0283 headroom, it carries
genuine dynamic range, and it is flat across timing. Had delayed exposure damaged the
underlying representation, that probe had the range to register it and did not.

The differential-susceptibility reading remains open, and its cost is specific rather than
unknown. Under the mean-endpoint readout it is unmeasurable at the \emph{n} achieved here. Under
a rate readout it is \textbf{un-run}: the exceedance machinery was only ever pointed at the
delay-0 question, and the seed inventory is asymmetric (24 delay-0 seeds per size, 3 delayed
seeds at 23M alone). A rate-collapse test at 60M --- 14 delayed seeds against the existing
24-seed baseline --- would settle it; at 23M the same test needs 82 and is out of reach (§8).

Finally, a note on what this result is not. Our pre-registration anticipated a
\emph{reopening}-side dissociation: some skill/intervention combinations recover and others do
not. What the data support is a \emph{detection}-side one --- syntax never showed a deficit, so no
syntax reopening was attempted. The two are superficially the same finding and are not the
same claim, and we report the one we have.

\section{§4 Reopen}\label{reopen}

\begin{figure}[t]
\centering
\rule{\linewidth}{36.00\baselineskip}
\caption{Placeholder for the §4 Reopen figure.}
\end{figure}

\subsection{4.1 An LR-schedule shift restores the delayed endpoint}\label{an-lr-schedule-shift-restores-the-delayed-endpoint}

All intervention arms fork from a single saved checkpoint at step 2500 --- identical weights,
identical optimizer state, identical step --- so that arms differ only in the intervention
applied. The \texttt{lr\_shift} arm restores the learning-rate schedule that a delay-0 model would
have been on at the point delayed exposure begins, then anneals normally. Success was
pre-registered as a two-part condition: the treated endpoint must reach 90 \% of that run's
\emph{own} early-exposure anchor, and general capability must not degrade beyond a matched
threshold.

\textbf{Table 3.} Endpoints at step 5000. The null is shown beside the treated arm throughout;
neither is collapsed into a ratio.

\begin{table}[t]
\centering\small
\begin{tabular}{@{}
  >{\raggedright\arraybackslash}p{(\linewidth - 12\tabcolsep) * \real{0.1429}}
  >{\raggedleft\arraybackslash}p{(\linewidth - 12\tabcolsep) * \real{0.1429}}
  >{\raggedleft\arraybackslash}p{(\linewidth - 12\tabcolsep) * \real{0.1429}}
  >{\raggedleft\arraybackslash}p{(\linewidth - 12\tabcolsep) * \real{0.1429}}
  >{\raggedleft\arraybackslash}p{(\linewidth - 12\tabcolsep) * \real{0.1429}}
  >{\raggedleft\arraybackslash}p{(\linewidth - 12\tabcolsep) * \real{0.1429}}
  >{\raggedleft\arraybackslash}p{(\linewidth - 12\tabcolsep) * \real{0.1429}}@{}}
\toprule\noalign{}
\begin{minipage}[b]{\linewidth}\raggedright
size
\end{minipage} & \begin{minipage}[b]{\linewidth}\raggedleft
treated (\texttt{lr\_shift})
\end{minipage} & \begin{minipage}[b]{\linewidth}\raggedleft
delayed null
\end{minipage} & \begin{minipage}[b]{\linewidth}\raggedleft
early anchor
\end{minipage} & \begin{minipage}[b]{\linewidth}\raggedleft
gate (0.90 × own d0)
\end{minipage} & \begin{minipage}[b]{\linewidth}\raggedleft
margin
\end{minipage} & \begin{minipage}[b]{\linewidth}\raggedleft
lift, \emph{p}
\end{minipage} \\
\midrule\noalign{}

23M & \textbf{0.767 ± 0.014} & 0.577 ± 0.051 & 0.791 ± 0.007 & 0.712 & +0.055 & 0.0065 \\
60M & \textbf{0.796 ± 0.022} & 0.600 ± 0.040 & 0.829 ± 0.027 & 0.746 & +0.050 & 0.0027 \\
\bottomrule
\end{tabular}
\end{table}


Per-seed treated endpoints are 0.752 / 0.778 / 0.772 at 23M and 0.796 / 0.818 / 0.774 at
60M. Every seed at both sizes clears its gate, and the per-seed lift over the matched null
survives Holm correction within the Reopen family (adjusted 0.0065 and 0.0054).

Two features of this table are more surprising than the recovery itself.

\textbf{The recovered endpoint is scale-invariant.} Both treated means fall inside the
pre-registered same-endpoint band {[}0.722, 0.812{]}, despite early anchors that differ by
0.038 and a 2.6× difference in parameters. Recovery converges on a level rather than
tracking capacity.

\textbf{The intervention collapses cross-seed variance.} The null's SD (0.051 at 23M, 0.040 at
60M) drops to the treated arm's (0.014, 0.022). The seed-variable severity documented in
§3.1 is a property of the \emph{untreated} delayed models; the intervention makes the outcome
both higher and more reproducible.

\subsection{4.2 The recovery does not come at the cost of general capability}\label{the-recovery-does-not-come-at-the-cost-of-general-capability}

Perplexity is reported as the three-number decomposition --- early anchor \emph{a}, delayed null
\emph{b}, treated arm \emph{c} --- because a single delta cannot distinguish a treated arm that
improved from a null that degraded.

At \textbf{23M there is no breach}: \emph{c} − \emph{b} = −0.889 / −0.204 / +0.193 across seeds, against a
pre-registered gate of +1.0 (seed 1: \emph{b} = 50.747, \emph{c} = 49.858, both from the canonical
\texttt{…seed1\_recheck} lineage). Two of three seeds improve.

At \textbf{60M the gate is tripped in all three seeds}: \emph{c} − \emph{b} = +1.670 ± 0.104 ppl. We report
it as a trip and honour it as written. Two controls bound what it can mean. A no-op fork
reproduces the null on both perplexity and accuracy (Δppl −0.012, Δacc +0.011), so fork
mechanics contribute nothing. And at \textbf{matched learning rate} the treated arm \emph{leads} the
null by 2.478 ± 0.240 ppl --- a lead larger than the breach --- because the arm is read
mid-anneal at roughly six times the null's terminal learning rate, with about 4.15 ppl of
annealing gain still unrealized. The breach persists only if the arm's remaining anneal
delivers less than 1.67 ppl of that. The evidence leans toward near-full erasure, but the
anneal slope is unmeasured, so we quote \textbf{no point estimate}: the residual cost is bounded
above by 1.67 ppl and a partly real cost is not excluded. Extending the arm past step 5000
is the run that would collapse this bound in either direction (§8).

\subsection{4.3 The effect is specific to the delayed skill}\label{the-effect-is-specific-to-the-delayed-skill}

The same intervention applied to delayed \emph{syntax} models --- replicated at all three seeds,
matching the depth of the arithmetic arm --- moves nothing: Δ local = −0.024 / −0.004 /
+0.044 and Δ LD = −0.021 / +0.007 / +0.015, sign-inconsistent around zero at every seed.
A single-seed control could not have distinguished ``no generic effect'' from ``an effect that
happened to land near zero once'', which is why the control was run at the arm's own seed
depth. Whatever \texttt{lr\_shift} does, it is not a generic late-training lift.

\subsection{4.4 Closure at this onset is largely a property of the schedule}\label{closure-at-this-onset-is-largely-a-property-of-the-schedule}

The remaining Tier-1 interventions were run at 23M seed 1 and are scored against the null's
cross-seed \textbf{range} {[}0.524, 0.626{]} rather than its mean, because a single seed landing
inside that range is indistinguishable from the null.

\begin{table}[t]
\centering\small
\begin{tabular}{@{}
  >{\raggedright\arraybackslash}p{(\linewidth - 4\tabcolsep) * \real{0.3333}}
  >{\raggedleft\arraybackslash}p{(\linewidth - 4\tabcolsep) * \real{0.3333}}
  >{\raggedright\arraybackslash}p{(\linewidth - 4\tabcolsep) * \real{0.3333}}@{}}
\toprule\noalign{}
\begin{minipage}[b]{\linewidth}\raggedright
arm (23M, seed 1)
\end{minipage} & \begin{minipage}[b]{\linewidth}\raggedleft
endpoint
\end{minipage} & \begin{minipage}[b]{\linewidth}\raggedright
verdict
\end{minipage} \\
\midrule\noalign{}

embedding reset & 0.586 & no detectable recovery (\emph{n} = 1) \\
reset + LR spike & 0.576 & no detectable recovery (\emph{n} = 1) \\
LR spike alone & 0.536 & no detectable recovery (\emph{n} = 1) \\
cram (concentrated re-exposure) & 0.574 & no detectable recovery (\emph{n} = 1) \\
\textbf{freeze} & \textbf{0.332} & \textbf{harmful} --- below the entire null range \\
\bottomrule
\end{tabular}
\end{table}


Only freeze supports a directional statement at this depth. The others are \emph{not} evidence
that those interventions fail; they are evidence that one seed cannot separate them from a
null whose own spread is 0.10 wide. Embedding reset was our pre-registered favourite, and
§7.2 states that prediction and its outcome as a standalone item; here it is one row of a
table, at \emph{n} = 1.

The one arm we extended to \emph{n} = 3 decides the mechanism. Concentrated re-exposure alone
leaves the endpoint at 0.574 --- at the null. Combined with the schedule shift it reaches
0.781 ± 0.031, far above cram alone and indistinguishable from the shift by itself: the
per-seed differences are −0.006 / +0.016 / +0.032, straddling zero with no consistent sign.
(We report the signs rather than the +0.014 mean, because a near-zero mean with disagreeing
signs is an absence of effect, not a small effect obscured by pooling.)

So additional skill density on top of a restored schedule adds nothing, while a restored
schedule without additional density recovers the skill. \textbf{At this onset, what closes the
window is largely the optimizer's decayed learning-rate schedule rather than representations
that have hardened against revision.} The scope of that statement is d50, arithmetic, and
these two sizes; it is not a claim that entrenchment plays no role at other onsets.

A final arm belongs in the record without supporting a claim. The 11M \texttt{lr\_shift} arm passes
its gate (0.747 against a gate of 0.484) but the pass is degenerate: the gate is 90 \% of an
anchor that failed the §3.1 floor, and the arm \emph{exceeds its own early anchor} (0.538),
which indicates that the 11M anchor was learning-rate-starved rather than that reopening
over-recovers. The 11M arm enters the paper only through its mechanism decomposition, where
the delayed lift of +0.389 splits into a generic component (+0.115) and an
exposure-alignment component (+0.274) --- about 70 \% of the effect requiring the delayed
exposure profile. That decomposition is exploratory and carries no headline.

\section{§5 Did the intervention repair the original circuit, or build a new one?}\label{did-the-intervention-repair-the-original-circuit-or-build-a-new-one}

Recovery tells us the skill came back; it does not tell us whether it came back \emph{in the same
machinery}. For each seed and size we compare a triad that shares an initialization: \textbf{A},
the early-exposure model; \textbf{B}, the delayed model; and \textbf{C}, the \texttt{lr\_shift}-recovered model
forked from B. Each model's arithmetic circuit is localized by mean-ablating attention heads
one at a time and recording exact-match damage. The decisive test is reciprocal: ablating
\textbf{A's} top heads inside \textbf{C} measures how much of A's circuit is still load-bearing in the
recovered model, and ablating \textbf{C's own} top heads inside C controls for C simply being
robust to any ablation. We report the transfer ratio --- the damage A's heads do in C as a
fraction of the damage they do in A --- where 1.0 means C inherited A's circuitry intact and 0
means A's heads are causally inert in C. One of the two size comparisons below is the
underpowered counterpart to §3.1's well-powered one; §5.3 makes that pairing explicit.

All numbers are computed on the single held-out problem set. An earlier version of this
analysis drew its eval items per-run, which overlapped training by 17--19 \% in four of the six
triads; all six were re-scored, the seed-1 triads reproduce exactly as a no-op control, and
the correction is reported in §2.4 because it changed a conclusion.

\subsection{5.1 Only one of four probes discriminates on this contrast}\label{only-one-of-four-probes-discriminates-on-this-contrast}

Before reading anything off the ablation numbers, it matters which methods can tell the two
outcomes apart \emph{at all}. Four were tried and three cannot:

\begin{itemize}
\tightlist
\item
  \textbf{Whole-layer mean-ablation saturates} --- ablating any core layer drives accuracy to
  approximately zero in every model, so it separates nothing. Localization is therefore done
  at head granularity.
\item
  \textbf{Linear CKA} of the residual stream sits at 0.66--0.83 across all six triads regardless of
  what the causal test says. Because the triads share an initialization, global geometry
  stays similar even where the wiring differs; CKA would have scored the most reorganized
  triad in the set at 0.658 and missed it entirely.
\item
  \textbf{Single-head activation patching} is flat: its restoration ratio spans 0.591--0.792
  (SD 0.083) against the ablation transfer's 0.187--0.868 (SD 0.265), and the two are
  uncorrelated (\emph{r} = 0.045, \emph{p} = 0.93; recomputing on the corrected transfers gives
  \emph{r} = −0.032, so the flatness is not an artifact of the contamination). The cause is a
  floor: single-head restoration is 0.05--0.15 in every model, so the ratio divides two small
  numbers and has no resolving power. An earlier draft cited patching as independent
  corroboration on the strength of one seed; running it on all six withdrew that.
\item
  \textbf{Pruning-based localization} was pre-screened and rejected: mask overlap spans only 0.088
  across the triads, and its agreement with ablation \emph{flips sign} across the sparsity sweep.
  A method whose sign is a hyperparameter choice cannot adjudicate (Appendix A).
\end{itemize}

We therefore report this cluster as \textbf{single-method by demonstration}, not by assumption.
Three non-discriminating probes is itself the evidence for that label --- the alternative,
citing agreement from a method never shown to resolve the contrast, is how the retracted
patching claim happened.

\subsection{5.2 Recovered circuits span a continuum, and it runs backwards at the shallow end}\label{recovered-circuits-span-a-continuum-and-it-runs-backwards-at-the-shallow-end}

\textbf{Table 4.} Reciprocal ablation transfer, ordered by deficit depth. Transfer = damage from
A's heads in C ÷ their damage in A. Control = damage from C's own heads ÷ damage from A's
heads in C, so a high control confirms C has a working circuit of its own.

\begin{table}[t]
\centering\small
\begin{tabular}{@{}
  >{\raggedright\arraybackslash}p{(\linewidth - 12\tabcolsep) * \real{0.1429}}
  >{\raggedleft\arraybackslash}p{(\linewidth - 12\tabcolsep) * \real{0.1429}}
  >{\raggedleft\arraybackslash}p{(\linewidth - 12\tabcolsep) * \real{0.1429}}
  >{\raggedleft\arraybackslash}p{(\linewidth - 12\tabcolsep) * \real{0.1429}}
  >{\raggedleft\arraybackslash}p{(\linewidth - 12\tabcolsep) * \real{0.1429}}
  >{\raggedleft\arraybackslash}p{(\linewidth - 12\tabcolsep) * \real{0.1429}}
  >{\raggedright\arraybackslash}p{(\linewidth - 12\tabcolsep) * \real{0.1429}}@{}}
\toprule\noalign{}
\begin{minipage}[b]{\linewidth}\raggedright
size
\end{minipage} & \begin{minipage}[b]{\linewidth}\raggedleft
seed
\end{minipage} & \begin{minipage}[b]{\linewidth}\raggedleft
deficit
\end{minipage} & \begin{minipage}[b]{\linewidth}\raggedleft
transfer
\end{minipage} & \begin{minipage}[b]{\linewidth}\raggedleft
control
\end{minipage} & \begin{minipage}[b]{\linewidth}\raggedleft
Jaccard(A, C)
\end{minipage} & \begin{minipage}[b]{\linewidth}\raggedright
pre-registered verdict
\end{minipage} \\
\midrule\noalign{}

23M & 2 & 0.182 & 0.659 & 1.50 & 0.00 & mixed / partial \\
60M & 2 & 0.202 & 0.404 & 2.43 & 0.00 & mixed / partial \\
23M & 3 & 0.210 & 0.868 & 1.16 & 0.20 & mixed / partial \\
60M & 1 & 0.224 & 0.445 & 2.18 & 0.09 & mixed / partial \\
60M & 3 & 0.260 & 0.201 & 5.13 & 0.00 & new circuit \\
23M & 1 & 0.292 & 0.187 & 5.36 & 0.00 & new circuit \\
\bottomrule
\end{tabular}
\end{table}


Recovered models do not behave alike. They span a continuum from near-complete transfer of
the original circuit to full reorganization. At one end (23M seed 3) the recovered model
inherits 87 \% of the damage A's circuit takes, and its own top heads are barely more damaging
than A's --- the original machinery is still carrying the skill. At the other (23M seed 1, 60M
seed 3) A's heads are close to inert in C while C's own heads are five times more damaging:
the skill returned in different machinery. \textbf{Every 60M triad sits at or below 0.445}, so at
the larger size reorganization is what we observe at every seed.

We report this axis continuously and assign no categorical verdict. Our implementation would
have labelled a triad \emph{repaired} only on a conjunction of high transfer and high
top-head-set overlap; the overlap condition is never met here --- the maximum is 0.20 --- so the
label was unreachable. It was an implementation default rather than a registered criterion,
and relabelling under a corrected rule would be choosing a decision rule after seeing the
number it decides (Appendix B). The mismatch is itself the methodological point: a recovered
model's head set can overlap its baseline's almost not at all while the baseline's heads
remain causally load-bearing --- near-zero overlap at transfer 0.868. \textbf{Set overlap and causal
overlap are different measurements}, and only the second answers the question. That is the
case for reciprocal ablation over any similarity-of-components measure, and the same failure
that makes CKA useless here, from a second direction.

The ordering suggests that deficit depth grades this outcome --- deeper deficits reorganize
more --- but the evidence does not survive its own robustness check. Across the six triads
Spearman ρ = −0.714 with an exact permutation \emph{p} = 0.0681, and leave-one-out is fragile:
four of the six single-point deletions push \emph{p} above 0.05. The direction persists and is not
resolvable at \emph{n} = 6. More triads at either size are the only thing that would settle it,
and this is the one place in the project where additional runs would still move a headline.

One pair of triads shows the relationship running backwards, and it is the best-measured
comparison in the table. At 23M, seeds 2 and 3 have nearly the same deficit (0.182 vs 0.210),
yet the \emph{deeper} one shows \emph{more} repair (0.868 vs 0.659), with non-overlapping bootstrap
intervals: {[}0.832, 0.903{]} against {[}0.611, 0.708{]}. Those intervals reflect eval-item variance
with one model per point --- they say the two measurements differ, not that two populations do
--- but within the shallow-deficit regime the severity ordering inverts. Whatever grades
circuit identity is not deficit depth alone.

\subsection{5.3 On size, this design is underpowered --- unlike the one in §3.1}\label{on-size-this-design-is-underpowered-unlike-the-one-in-3.1}

Transfer averages 0.571 ± 0.349 at 23M and 0.350 ± 0.131 at 60M (Welch \emph{p} = 0.39). \textbf{That
non-significance means nothing here.} The minimum detectable effect at 80 \% power is 0.809,
against an outcome span of only 0.681 across all six triads: no difference this design could
have produced would have been detected. The question is \emph{untested}, not answered --- and it is
no longer unidentifiable, since the corrected deficit ranges at the two sizes overlap, so
size is separable from deficit depth in a way it was not before the contamination fix.

This is worth setting beside §3.1's size comparison, which is structurally identical --- same
\emph{n} = 3 per group, same Welch test, same two sizes, also non-significant --- and carries
opposite weight. There, the MDE (0.144) sits below the deficit being asked about (0.21), so
the null genuinely excludes a large size effect. Here it exceeds the entire measured range.
The two together are the clearest illustration of a rule this paper applies throughout:
non-significance is read off the minimum detectable effect, never off the \emph{p}.

What the size axis does show is \textbf{variance}. The 60M transfers are tightly clustered
(SD 0.131) where the 23M ones are not (SD 0.349), a 2.7× ratio, mirroring §4's finding that
the intervention collapses cross-seed variance in the endpoint. We report this as suggestive:
it is a single comparison at \emph{n} = 3 per group on an observational rather than manipulated
factor.

\subsection{5.4 What these results do not speak to}\label{what-these-results-do-not-speak-to}

Our third pre-registered prediction linked circuit identity to the reset intervention: if
reset were the winning intervention, we expected it to build a new circuit. Reset never
produced recovery (§4.4), so the antecedent never occurred and the prediction was never
entered. The models diagnosed here are \texttt{lr\_shift}-recovered, which that prediction did not
cover; §7.2 states this rather than letting the section read as an answer to it.

Three further quantities appear in our released logs and support nothing. The within-60M
Spearman is mechanical at \emph{n} = 3. An extrapolation from the 23M triads onto the 60M points
was built before the contamination correction and no longer fits (residuals −0.30, −0.15,
−0.20). Numbers from the causal-tracing and pruning scripts were computed on the contaminated
draws and were never re-run, since neither supports a surviving claim. Circuit redundancy and
sufficiency are likewise untested here: the retracted patching analysis was the only probe
that addressed them.

\section{Related Work}\label{related-work}

\textbf{Critical periods.} The concept comes from development, where early sensory deprivation
produces deficits that later normal input does not repair \citep{hubel-wiesel-1970,lenneberg-1967,johnson-newport-1989}. The closest machine-learning analogue is Achille et al.'s demonstration
that deep networks have critical learning periods: a deficit imposed early in training
(blurred images) leaves a permanent gap even given unlimited clean data afterwards, and the
window's timing is a property of training rather than of architecture \citep{achille-etal-2019}. That work manipulates \emph{input quality} in vision models; we manipulate
\emph{when a skill's data arrives} in a language model while holding its total quantity fixed,
which is what lets a deficit be attributed to timing rather than to dose or degraded signal.

\textbf{Plasticity loss and warm starting.} A separate literature finds that networks become
harder to train as training proceeds: warm-started models generalize worse than freshly
initialized ones given the same final data \citep{ash-adams-2020}, and continual-learning
agents progressively lose the ability to fit new tasks, recoverable by re-injecting
plasticity through resets or noise \citep{dohare-etal-2024,lyle-etal-2023,berariu-etal-2021}.
\textbf{This is the closest prior explanation of our reopening result and the one we most risk
rediscovering}, so we state the difference precisely: those settings involve a distribution
shift, a restart, or a sequence of distinct tasks, and typically recover plasticity by
perturbing \emph{weights}. Our models train once, without interruption, on a stationary mixture
whose only manipulated variable is a skill's onset step --- and the intervention that recovers
the skill changes no weights at all, only the learning-rate schedule, while resetting weights
(the standard plasticity remedy) produced no detectable recovery in our setting at \emph{n} = 1.
The mechanisms may still be the same one seen from two directions; what we can say is that
the recovery here is obtainable without weight perturbation. (That reset arm is \emph{n} = 1;
§7.2 states what it does and does not license.)

\textbf{Curriculum, data ordering, and delayed generalization.} Ordering effects have a long
history \citep{bengio-etal-2009} with mixed results at LM scale, where reordering often yields
little; our finding is compatible with that, since total exposure is what matters for most
of the corpus and it is one late-arriving skill that pays a timing cost. Delayed
generalization on arithmetic is also studied as grokking \citep{power-etal-2022}, which is a
different phenomenon despite the surface adjacency: grokking is a long plateau followed by
sudden generalization at \emph{fixed} data, whereas ours is a permanent endpoint gap produced by
\emph{shifting} when the data arrives.

\textbf{Small-scale LM evaluation.} We use the BabyLM setting \citep{warstadt-etal-2023} and
BLiMP-style minimal pairs \citep{warstadt-etal-2020}, with the metric modification of §2.3, and
the agreement-attraction literature \citep{linzen-etal-2016,marvin-linzen-2018}
for §6.

\textbf{Circuit identity.} Our diagnosis uses causal ablation and activation patching in the
circuit-analysis tradition \citep{elhage-etal-2021,wang-etal-2023,meng-etal-2022} and
representational similarity \citep{kornblith-etal-2019}; §5.1 reports the negative
methodological result that three of those four probes do not discriminate on our contrast.

\begin{center}\rule{0.5\linewidth}{0.5pt}\end{center}

\section{§6 A secondary probe: agreement attraction across scale}\label{a-secondary-probe-agreement-attraction-across-scale}

The models trained for this study support one question the critical-period design did not
set out to ask. Subject--verb agreement is acquired almost perfectly here --- 0.98 on
controlled items --- yet fails when a number-mismatching noun intervenes. A fully-crossed item
design, in which the verb-form frequency prior cancels \emph{exactly},
separates two components of that failure: \textbf{attraction}, the cost of an intervening noun
whose number mismatches the head, and \textbf{complexity}, the cost of any intervening material
at all.

\textbf{Attraction is real, large, and the dominant component.} At 23.2M it costs
+0.303 ± 0.064 (\emph{p} = 8.6e-05) against complexity's +0.143 ± 0.035, in both number
directions and both frame types. A unigram-frequency control on the same items returns
exactly 0.000 by construction, so the effect cannot be a verb-form prior. Attraction exceeds
complexity at every size and is positive in \textbf{81 of 81 models} across two independent
training mixtures --- 72 syntax-enriched (3 sizes × 24 seeds) and 9 arithmetic-enriched.

\textbf{Whether these models are anti-syntactic is scale-dependent.} Mismatch accuracy on
controlled items is 0.405 at 11.5M (below chance, \emph{p} \textless{} 0.001), 0.511 at 23.2M (at chance,
\emph{p} = 0.435), and 0.610 at 60.3M (above chance, \emph{p} \textless{} 0.001). Below-chance agreement occurs
only at the smallest size.

\textbf{Attraction declines with scale; whether it dissociates from complexity depends on the
mixture.} In the syntax-enriched mixture attraction falls \textbf{−0.316 per decade of
parameters} (\emph{r} = −0.803, \emph{p} = 2.3e-17, \emph{n} = 72), from +0.430 to +0.201, while complexity
is \textbf{flat} (+0.016/decade, \emph{p} = 0.43). The second mixture reproduces the attraction half
in direction and magnitude, \textbf{−0.207/decade}, stable at −0.21 to −0.24 under every
estimator tried but at \textbf{estimator-dependent significance --- \emph{p} = 0.033 under OLS, 0.073
under a heteroscedasticity-robust SE} --- on an unbalanced (3, 3, 8) seed design whose
leverage sits on the under-sampled small end. It does \textbf{not} reproduce the dissociation:
complexity falls there too, and the paired within-model contrast is flat (−0.095,
\emph{p} = 0.36) against −0.333 (\emph{p} = 3.8e-11) in the syntax mixture, the one part of the
comparison that is not estimator-fragile. A cross-mixture divergence at 60.3M \textbf{does not
survive replication} at \emph{n} = 8 (+0.065, \emph{p} = 0.091, halved from +0.133) and \textbf{is not
excluded} --- the 95 \% CI still contains the original effect. An English-exposure explanation
for it is excluded separately: exposure is constant across sizes by design and can only
produce a uniform offset.

\textbf{Within a single run, attraction may also decline over training, but we do not claim it.}
Between steps 2500 and 5000 the isolated attraction cost falls +0.392 → +0.303 (paired,
\emph{n} = 6 at 23.2M, 5/6 seeds) --- marginal at \emph{p} = 0.053, with 53 \% power against this effect
size and an MDE of 0.125; roughly ten seeds would decide it. The mismatch cell itself
improves significantly over the same interval (+0.070, \emph{p} = 0.041). We report the direction
and the power, and treat the trajectory as suggestive only.

\textbf{The standard metric gets this wrong at every scale, in three different ways.} Uncontrolled
BLiMP-style items report below-chance behaviour at all three sizes, which is true only at the
smallest; the frequency-controlled metric locates it correctly; and only the crossed design
separates attraction from complexity. Three metrics, three conclusions, on the same models.
A related trap is worth naming: pooling all 72 models into one attraction estimate averages
across a size-dependent effect and yields a number that describes no size --- an
averaging-across-a-dissociation error, not a confound.

This is also why the skill sits here rather than in §3. Agreement is unusable as a
delayed-exposure target: it saturates by step 2500 and its supporting structure appears in
roughly 39 \% of the corpus against \emph{wh}-extraction's 1.5 \%, so it cannot meaningfully be
withheld. Being \emph{learnable enough to measure} and being \emph{delayable} are different properties,
and a design like ours needs both.

\section{§7 Discussion}\label{discussion}

\subsection{7.1 Two timing results at different levels}\label{two-timing-results-at-different-levels}

This paper reports two findings about timing, and they operate at different levels. The
first says that \emph{which} skills a late start damages is skill-specific: arithmetic pays a
persistent cost for a delayed start, long-distance \emph{wh}-licensing pays nothing we can
measure. The second says that where the damage does occur, it is carried by the optimizer's
learning-rate schedule rather than by representations that have hardened against revision ---
restoring the schedule restores the skill, and adding skill data without restoring the
schedule does not. A late start hurts some skills and not others; when it hurts, what has
closed is the model's step size, not its mind.

\subsection{7.2 The three pre-registered predictions}\label{the-three-pre-registered-predictions}

\textbf{Prediction \#1 --- that syntax would show the \emph{stronger} critical period --- was inverted.}
Arithmetic carries the effect at both confirmatory sizes; long-distance \emph{wh}-licensing shows
no detectable deficit under five distinct probes. We had reasoned that long-distance
dependencies, being scaffolded by earlier-acquired structure, would be the more fragile to a
late start. The opposite skill turned out to be the fragile one, and this inversion is the
origin of the dissociation in §3.3 rather than an incidental miss.

\textbf{Prediction \#2 --- that embedding reset would be the most effective intervention --- was not
borne out.} Reset produced no detectable recovery, landing at 0.586 against a null whose
cross-seed range is {[}0.524, 0.626{]}; the learning-rate shift was the only arm to clear the
pre-registered gate. Because reset was run at a single seed, this is a statement about what
we could detect, not a demonstration that reset fails: one seed cannot be separated from a
null that spans 0.10. The prediction's second half --- that reset would carry the largest
forgetting risk --- was also not observed as stated.

\textbf{Prediction \#3 --- that a reset-driven recovery would build new circuitry --- was never
testable, because its antecedent did not occur.} The prediction was conditional on reset
being the winning intervention. Reset never produced recovery, so the conditional was never
entered. §5 diagnoses models recovered by the learning-rate shift, which prediction \#3 did
not cover; those results should not be read as confirming or disconfirming it. We state this
explicitly because §5 occupies the position in the paper where a reader holding our
pre-registration will look for \#3's answer.

\subsection{7.3 What the mechanism result does not say}\label{what-the-mechanism-result-does-not-say}

Schedule-driven closure is demonstrated at one onset (the midpoint of training), for one
skill, at two model sizes. It is not a claim that representational entrenchment plays no role
in critical periods generally, nor that it plays no role at later onsets in this same setup.
Our intermediate-onset runs exist at a single seed and are not evidence about the shape of
the timing curve.

\begin{center}\rule{0.5\linewidth}{0.5pt}\end{center}

\section{§8 Limitations}\label{limitations}

Five specific experiments would move a label, and we name them with their costs rather than
gesturing at future work.

\textbf{The dissociation's stronger reading is un-run, not unknowable.} Whether syntax is
genuinely \emph{less susceptible} to delayed exposure, rather than merely having too little
headroom to show susceptibility, is decidable by a rate-collapse test at 60M: 14 delayed
seeds against our existing 24-seed baseline, or 6 at a permissive threshold. This assumes
delay eliminates the high-acquisition mode; if it merely halves it, 59 are needed. The same
test at 23M requires 82 seeds and is out of reach. This is the cheapest outstanding
experiment in the project and the only one that can promote the dissociation past its
current scope.

\textbf{The 60M perplexity cost is bounded, not estimated.} Extending the treated arm past step
5,000 to complete its annealing would convert our upper bound of 1.67 ppl into a point
estimate in either direction.

\textbf{The circuit-identity gradient is unresolved at \emph{n} = 6.} Additional triads at either size
are the only thing that would settle whether deficit depth grades circuit identity; the
direction persists but four of six leave-one-out deletions break its significance.

\textbf{The circuit results are single-method by demonstration.} Joint head-group patching, shown
to have dynamic range on the repair-versus-reorganization contrast, would let those claims
drop the annotation. Three probes we tried do not discriminate.

\textbf{One agreement-attraction slope is estimator-dependent.} Three or more further seeds at the
smallest size would rebalance the leverage that makes it significant under OLS and marginal
under a robust standard error.

Beyond these: one corpus, one tokenizer, one language, a maximum of 60.3M parameters, and a
single delay onset carrying the reopening result. Five of six intervention arms are single
seed. The smallest size is excluded by a floor rather than solved.

\begin{center}\rule{0.5\linewidth}{0.5pt}\end{center}

\section{§9 Conclusion}\label{conclusion}

Timing of exposure matters in these models, but not uniformly and not for the reason we
expected. Delayed arithmetic exposure leaves a deficit that further training does not close;
delayed long-distance \emph{wh}-licensing leaves nothing we can detect, which inverts our
prediction and which we state in absolute terms because the two skills are not equally
observable. Where the deficit exists it reopens, and it reopens when the learning-rate
schedule is restored rather than when the missed data is concentrated --- so what closes at
this onset behaves like a property of the optimizer rather than of entrenched representation.
Whether the skill contrast reflects genuine differential susceptibility remains open, and we
have named the fourteen-seed experiment that would decide it.

\begin{center}\rule{0.5\linewidth}{0.5pt}\end{center}

\end{document}
